\documentclass{article}
\usepackage{amsmath,amssymb}
\usepackage[vocab={diophantine}]{leanify}
\newtheorem{lemma}{Lemma}
\begin{document}

\section{Notation}

We write $g(a,b)$ for the Frobenius number of $a$ and $b$.
We write $\gamma(a,b)$ for the genus of the numerical semigroup generated by $a$ and $b$.
We write $N(n,a,b)$ for the number of representations of $n$ by $a$ and $b$.
We write $r_2(n)$ for the number of representations of $n$ as a sum of two squares.
We write $X(m,n)$ for the first leg of the Pythagorean triple generated by $m$ and $n$.
We write $Y(m,n)$ for the second leg of the Pythagorean triple generated by $m$ and $n$.
We write $Z(m,n)$ for the hypotenuse of the Pythagorean triple generated by $m$ and $n$.

\section{Linear equations}

\begin{lemma}\label{sylvester}
$g(3,5)=7$, $\gamma(3,5)=4$, $N(7,3,5)=0$ and $N(8,3,5)=1$.
\end{lemma}

\section{Quadratic equations}

\begin{lemma}\label{squares}
$r_2(5)=8$ and $r_2(3)=0$.
\end{lemma}

\begin{lemma}\label{euclid}
Let $m$ and $n$ be positive integers with $n<m$. Then $X(m,n)^2+Y(m,n)^2=Z(m,n)^2$.
\end{lemma}

\begin{lemma}\label{congruent}
$6$ is a congruent number and $13$ is representable as a sum of two squares.
\end{lemma}

\begin{lemma}\label{representable}
Let $n$ be an integer with $n\ge0$ and $r_2(n)>0$. Then $n$ is representable as a sum of two squares.
\end{lemma}

\end{document}
